\begin{trapbox}
	\textbf{\textcolor{CTrap}{易错点一：只看两段，漏看接口}}

	两段都递增，整体也可能不递增。例如
	$f(x)=\begin{cases}x+2,&x\le0,\\x,&x>0.\end{cases}$
	左端值为 $2$，右极限为 $0$；从 $0$ 往右发生向下跳跃。递增须另查 $g(x_0)\le L$。

	\medskip
	\textbf{\textcolor{CTrap}{易错点二：把接口不等号误写成严格不等号}}

	虽然要求\emph{严格}递增，接口仍是 $g(x_0)\le L$。
	例如 $f(x)=x$ 可写成 $g(x)=x\ (x\le0)$、$h(x)=x\ (x>0)$；此时 $g(0)=L=0$，但 $f$ 在全域严格递增。

	\medskip
	\textbf{\textcolor{CTrap}{易错点三：递减时方向写反}}

	沿着 $x$ 增大的方向往右看，整体递减须“左高右低”，即 $g(x_0)\ge L$。

	\medskip
	\textbf{\textcolor{CTrap}{易错点四：直接把 $h(x_0)$ 当作右段函数值}}

	右段条件为 $x>x_0$，接口属于左段，$f(x_0)=g(x_0)$。应比较 $g(x_0)$ 与右极限 $L$；即使表达式 $h$ 在 $x_0$ 可代入，也只把代入结果用于\emph{求极限}。

	\medskip
	\textbf{\textcolor{CTrap}{易错点五：以为接口必须连续}}

	向单调方向的跳跃允许存在。例如 $f(x)=\begin{cases}x,&x\le0,\\x+1,&x>0\end{cases}$ 在全域严格递增，接口却不连续。
\end{trapbox}